% ============================================================
% Section 172: 中心势垒微扰下的基态能量
% ============================================================
\section{量子力学：中心势垒微扰下的基态能量}
\label{sec:center-barrier-perturbation}

\begin{modelbox}[题目：中心势垒微扰下的基态能量]
质量为 $m$ 的粒子在一维势阱中运动，势函数为
\[
V(x)=\begin{cases}\infty, & |x|>3a \\ 0, & a<|x|<3a \\ V_0, & |x|<a\end{cases}
\]
将 $V_0$ 部分视为在宽度为 $6a$ 的平坦无限深势阱上的微扰，用一级微扰方法计算基态能量。
\end{modelbox}

\begin{tipbox}[考点快速分析]
\begin{itemize}
    \item \textbf{未微扰系统}：宽度 $6a$ 的无限深方势阱，基态为偶宇称余弦波
    \item \textbf{微扰项}：中心区域 $|x|<a$ 的势垒 $V_0$
    \item \textbf{一级修正}：$E_1^{(1)}=\langle\psi_1^{(0)}|H'|\psi_1^{(0)}\rangle=V_0\int_{-a}^{a}|\psi_1|^2dx$
    \item \textbf{积分技巧}：利用 $\cos^2\theta=(1+\cos2\theta)/2$
\end{itemize}
\end{tipbox}

\begin{derivationbox}[标准解题过程]
\textbf{1. 未微扰系统}

宽度 $L=6a$ 的无限深方势阱：
\[
\psi_n^{(0)}(x)=\sqrt{\frac{2}{L}}\sin\frac{n\pi(x+3a)}{L}=\sqrt{\frac{1}{3a}}\sin\frac{n\pi(x+3a)}{6a},
\]
能级 $E_n^{(0)}=\dfrac{n^2\pi^2\hbar^2}{72ma^2}$。

对基态 $n=1$（偶宇称），用更简洁的形式：
\[
\psi_1^{(0)}(x)=\sqrt{\frac{1}{3a}}\cos\frac{\pi x}{6a},\quad E_1^{(0)}=\frac{\pi^2\hbar^2}{72ma^2}.
\]

\textbf{2. 一级微扰能量修正}

微扰势 $H'(x)=V_0$（$|x|<a$），其余为零。一级修正：
\[
E_1^{(1)}=\int_{-a}^{a}V_0\,|\psi_1^{(0)}(x)|^2\,dx=\frac{V_0}{3a}\int_{-a}^{a}\cos^2\frac{\pi x}{6a}\,dx.
\]

利用 $\cos^2\theta=\frac{1+\cos2\theta}{2}$：
\begin{align*}
\int_{-a}^{a}\cos^2\frac{\pi x}{6a}\,dx&=\int_{-a}^{a}\frac12\Bigl[1+\cos\frac{\pi x}{3a}\Bigr]dx\\
&=a+\frac{3a}{2\pi}\Bigl[\sin\frac{\pi}{3}-\sin\Bigl(-\frac{\pi}{3}\Bigr)\Bigr]\\
&=a+\frac{3a}{2\pi}\cdot\sqrt3=a\Bigl(1+\frac{3\sqrt3}{2\pi}\Bigr).
\end{align*}

因此
\begin{equation}
\boxed{E_1^{(1)}=\frac{V_0}{3a}\cdot a\Bigl(1+\frac{3\sqrt3}{2\pi}\Bigr)=V_0\Bigl(\frac13+\frac{\sqrt3}{2\pi}\Bigr).}
\end{equation}

\textbf{3. 微扰后的基态能量}

\begin{equation}
\boxed{E_1\approx E_1^{(0)}+E_1^{(1)}=\frac{\pi^2\hbar^2}{72ma^2}+V_0\Bigl(\frac13+\frac{\sqrt3}{2\pi}\Bigr).}
\end{equation}
\end{derivationbox}

\begin{formulabox}[答案汇总]
\begin{center}
\begin{tabular}{ll}
\toprule
\textbf{结论} & \textbf{结果} \\
\midrule
未微扰基态能量 & $E_1^{(0)}=\pi^2\hbar^2/(72ma^2)$ \\
一级能量修正 & $E_1^{(1)}=V_0\bigl(\frac13+\frac{\sqrt3}{2\pi}\bigr)$ \\
微扰后基态能量 & $E_1=E_1^{(0)}+E_1^{(1)}$ \\
\bottomrule
\end{tabular}
\end{center}
\end{formulabox}

\begin{insightbox}[物理图像]
\begin{itemize}
    \item 中心势垒对基态能量的修正正比于 $V_0$ 和波函数在势垒区域的占据概率。因基态波函数在中心有最大值（波腹），修正量相对较大。
    \item 若势垒位于波节位置（如 $n=2$ 的 $x=0$ 处），一级修正将为零。本题中 $n=1$ 的基态在 $x=0$ 处概率密度最大，故 $E_1^{(1)}$ 显著非零。
    \item 当 $V_0\to\infty$ 时，微扰论失效，体系退化为两个独立的宽度 $2a$ 势阱（中间硬壁），基态能量应为 $\pi^2\hbar^2/(8ma^2)$，远大于原宽度 $6a$ 的基态。
\end{itemize}
\end{insightbox}
